\documentclass[12pt,a4paper]{article}
\usepackage[a4paper,margin=25mm,headheight=15pt]{geometry}
\usepackage[T1]{fontenc}
\usepackage{lmodern}
\usepackage{xcolor}
\usepackage{fancyhdr}
\usepackage{setspace}
\usepackage{sectsty}
\usepackage{enumitem}
\usepackage{hyperref}
\hypersetup{colorlinks=true,linkcolor=blue,urlcolor=blue}
\pagestyle{fancy}
\fancyhf{}
\lhead{\textit{LaTeX Learning Module}}
\rhead{\textit{Page Formatting}}
\cfoot{\thepage}
\renewcommand{\headrulewidth}{0.4pt}
\renewcommand{\footrulewidth}{0pt}
\sectionfont{\color{blue}}
\subsectionfont{\color{teal}}
\setlist[itemize]{itemsep=3pt,topsep=5pt}
\title{\textbf{Page Layout and Document Styling}}
\author{LaTeX Learning Module}
\date{\today}
\begin{document}
\maketitle
\thispagestyle{plain}
\section{Choosing a Page Style}
The \verb|\pagestyle| command controls the header and footer style used on pages throughout a document. LaTeX includes several built-in choices:
\begin{itemize}
\item \texttt{plain} places the page number at the bottom of the page.
\item \texttt{empty} removes the standard header and footer.
\item \texttt{headings} displays running headings and page numbers.
\end{itemize}
Use \verb|\thispagestyle{empty}| when a change should affect only the current page. Use \verb|\pagestyle{empty}| when it should apply to the whole document.
\section{Setting Margins with geometry}
The \texttt{geometry} package makes it straightforward to set page dimensions. For example:
\begin{verbatim}
\usepackage[a4paper,margin=25mm]{geometry}
\end{verbatim}
You can also set each edge separately:
\begin{verbatim}
\usepackage[
  top=25mm,
  bottom=25mm,
  left=30mm,
  right=20mm
]{geometry}
\end{verbatim}
Choose margins that leave enough room for the text and any printed or binding requirements.
\section{Creating Headers and Footers}
The \texttt{fancyhdr} package allows different information to appear in the left, center, and right areas of the header and footer.
\begin{verbatim}
\usepackage{fancyhdr}
\pagestyle{fancy}
\fancyhf{}
\lhead{Course Notes}
\rhead{Page Layout}
\cfoot{\thepage}
\end{verbatim}
The commands \verb|\lhead|, \verb|\chead|, and \verb|\rhead| set the header areas. Their footer equivalents are \verb|\lfoot|, \verb|\cfoot|, and \verb|\rfoot|. The command \verb|\fancyhf{}| clears existing header and footer fields before new ones are assigned.
\section{Paragraph Indentation and Spacing}
LaTeX indents the first line of many paragraphs by default. Set the indentation length in the preamble with \verb|\parindent|:
\begin{verbatim}
\setlength{\parindent}{0pt}
\end{verbatim}
The distance between paragraphs can be adjusted with \verb|\parskip|:
\begin{verbatim}
\setlength{\parskip}{6pt}
\end{verbatim}
Use these settings consistently. Large paragraph gaps combined with first-line indentation may make a page look uneven.
\section{Line Spacing}
The \texttt{setspace} package provides convenient commands for changing line spacing:
\begin{verbatim}
\usepackage{setspace}
\singlespacing
\onehalfspacing
\doublespacing
\end{verbatim}
You can apply a spacing choice to the whole document or use the \texttt{spacing} environment for a selected portion:
\begin{verbatim}
\begin{spacing}{1.25}
This paragraph uses one-and-a-quarter spacing.
\end{spacing}
\end{verbatim}
\onehalfspacing
This paragraph is included as a visual example of the wider line spacing enabled by \verb|\onehalfspacing|. It can be useful for drafts that need room for comments or review.
\singlespacing
\section{Horizontal and Vertical Space}
Use \verb|\hspace| to add horizontal space and \verb|\vspace| to add vertical space.
\begin{verbatim}
Text\hspace{8mm}More text
\vspace{5mm}
A new line after some vertical space.
\end{verbatim}
Manual spacing is useful for small adjustments, but avoid using it as the main method for arranging an entire page. LaTeX's layout tools usually produce more reliable results.
\section{Page Breaks}
The commands \verb|\newpage| and \verb|\pagebreak| both request a new page. They differ in how strongly they encourage the break:
\begin{itemize}
\item \verb|\newpage| ends the current page immediately.
\item \verb|\pagebreak| suggests a page break and can take an optional strength value from 0 to 4.
\end{itemize}
For example, \verb|\pagebreak[3]| requests a strong break without making it quite as absolute as \verb|\newpage|.
\section{Customizing Headings, Colors, and Lists}
The \texttt{sectsty} package can change the appearance of section headings. For example:
\begin{verbatim}
\usepackage{sectsty}
\sectionfont{\color{blue}}
\subsectionfont{\color{teal}}
\end{verbatim}
The \texttt{xcolor} package provides named colors and lets you define your own. The \verb|\pagecolor| command changes the background color of a page:
\begin{verbatim}
\usepackage{xcolor}
\pagecolor{yellow!8}
\end{verbatim}
Use a light background when the document is intended for reading or printing. To return to the normal background, use \verb|\pagecolor{white}|.
The \texttt{enumitem} package gives more control over lists. For example:
\begin{verbatim}
\usepackage{enumitem}
\setlist[itemize]{itemsep=4pt,leftmargin=18mm}
\end{verbatim}
You can also change the settings for one list without changing every list in the document:
\begin{verbatim}
\begin{itemize}[itemsep=6pt]
  \item First point
  \item Second point
\end{itemize}
\end{verbatim}
\section{Keeping Text Together with mbox}
The \verb|\mbox| command keeps its contents on one line, preventing LaTeX from breaking the text at spaces inside the box.
\begin{verbatim}
The product code is \mbox{AB 2048 Z}.
\end{verbatim}
If the complete box cannot fit on the line, it may extend into the margin. Use this command only for short content that should remain together.
\section{Avoiding Page Splits with samepage}
The \texttt{samepage} environment asks LaTeX to keep a group of material on one page rather than splitting it across a page boundary.
\begin{verbatim}
\begin{samepage}
A short heading
A paragraph that should stay with the heading.
\end{samepage}
\end{verbatim}
This is useful for compact groups of related content. It does not guarantee that an oversized block will fit on a page.
\section{Creating Text Blocks with parbox}
The \verb|\parbox| command places paragraph text inside a box with a chosen width. LaTeX wraps the text to fit that width.
\begin{verbatim}
\parbox{0.7\textwidth}{
  This text is placed inside a box that uses
  seventy percent of the current text width.
}
\end{verbatim}
A \verb|\parbox| can be useful for short notes, side-by-side text areas, or small formatted blocks.
\section{Adjusting Vertical Position with raisebox}
The \verb|\raisebox| command moves content up or down relative to the surrounding text baseline.
\begin{verbatim}
Normal text \raisebox{2pt}{raised text}
and \raisebox{-2pt}{lowered text}.
\end{verbatim}
Positive values raise the content; negative values lower it. Small adjustments are usually easiest to read.
\section{Manual Line Breaks}
Use \verb|\\| to end a line where a line break is appropriate, such as in an address or a short title block. The \verb|\newline| command can also create a line break.
\begin{verbatim}
North Street 12\\
Central District\\
Example City
\end{verbatim}
For ordinary paragraphs, leave a blank line in the source to start a new paragraph instead of inserting repeated manual line breaks.
\section{Key Points}
Page styles control running information, while page-layout settings determine the space available for content. Packages such as \texttt{fancyhdr}, \texttt{sectsty}, and \texttt{enumitem} provide additional control over the document's appearance. Text-positioning commands are best used selectively, so the layout remains clear and consistent.
\end{document}
